\documentclass[11pt,a4paper]{article}
\usepackage[utf8]{vietnam}
\usepackage{amsmath,amssymb}
\usepackage{tabularx}
\usepackage{tikz}
\usepackage{xcolor}
\usepackage{fancyhdr}
\usepackage[top=1.8cm, bottom=2.0cm, left=1.8cm, right=1.8cm]{geometry}

\setlength{\parskip}{3pt}
\setlength{\parindent}{0pt}

\pagestyle{fancy}
\fancyhf{}
\renewcommand{\headrulewidth}{0.4pt}
\lhead{\footnotesize \textbf{CHÙA HƯƠNG SƠN -- HƯỚNG DẪN CHẤM THI}}
\rhead{\footnotesize \textbf{ĐÁP ÁN CHI TIẾT ĐỀ THI THỬ SỐ 01}}
\cfoot{\thepage}

\begin{document}

\begin{center}
{\large\bfseries ĐÁP ÁN \& HƯỚNG DẪN CHẤM ĐỀ THI THỬ CUỐI HỌC KỲ I (ĐỀ SỐ 01)}\\[1.5mm]
{\bfseries MÔN: TOÁN -- LỚP 9 \quad|\quad NĂM HỌC 2026 -- 2027}\\[1mm]
\textit{(Chuẩn cấu trúc định dạng mới của Sở Giáo dục và Đào tạo)}
\end{center}

\vspace{1mm}
\hrule
\vspace{3mm}

\noindent
\textbf{\large PHẦN I. TRẮC NGHIỆM KHÁCH QUAN (7,0 ĐIỂM)}

\vspace{2mm}
\noindent
\textbf{1. Bảng đáp án Phần 1: Trắc nghiệm nhiều lựa chọn (3,0 điểm)} \textit{(Mỗi câu đúng được 0,25 điểm)}
\begin{center}
\begin{tabular}{|c|c|c|c|c|c|c|c|c|c|c|c|c|}
\hline
\textbf{Câu} & \textbf{1} & \textbf{2} & \textbf{3} & \textbf{4} & \textbf{5} & \textbf{6} & \textbf{7} & \textbf{8} & \textbf{9} & \textbf{10} & \textbf{11} & \textbf{12} \\
\hline
\textbf{Đáp án} & \textbf{A} & \textbf{A} & \textbf{C} & \textbf{B} & \textbf{B} & \textbf{C} & \textbf{A} & \textbf{B} & \textbf{A} & \textbf{B} & \textbf{C} & \textbf{C} \\
\hline
\end{tabular}
\end{center}

\vspace{2mm}
\noindent
\textbf{2. Bảng đáp án Phần 2: Trắc nghiệm Đúng -- Sai (2,0 điểm)}
\begin{center}
\begin{tabular}{|c|c|c|c|c|l|}
\hline
\textbf{Câu} & \textbf{Ý a} & \textbf{Ý b} & \textbf{Ý c} & \textbf{Ý d} & \textbf{Ghi chú đáp án chuẩn} \\
\hline
\textbf{Câu 13} & \textbf{Đ} & \textbf{Đ} & \textbf{S} & \textbf{Đ} & Ý c) sai do vế phải là $820$ chứ không phải $800$. \\
\hline
\textbf{Câu 14} & \textbf{Đ} & \textbf{Đ} & \textbf{S} & \textbf{Đ} & Ý c) sai do chu vi $C = 2\pi R = 10\pi\text{ cm}$ (đề ghi $25\pi$). \\
\hline
\end{tabular}
\end{center}
\textit{Quy tắc tính điểm Bộ GD\&ĐT:} Đúng 1 ý: 0,1đ; Đúng 2 ý: 0,25đ; Đúng 3 ý: 0,5đ; Đúng cả 4 ý: 1,0đ.

\vspace{2mm}
\noindent
\textbf{3. Bảng đáp án Phần 3: Trắc nghiệm trả lời ngắn (2,0 điểm)} \textit{(Mỗi câu đúng được 0,5 điểm)}
\begin{center}
\begin{tabular}{|c|l|c|c|}
\hline
\textbf{Câu} & \textbf{Nội dung câu hỏi} & \textbf{Kết quả bài làm} & \textbf{Điểm} \\
\hline
\textbf{Câu 15} & Nghiệm của bất phương trình $3(x + 2) - 5 > 2(2x - 1) + 1$ & \textbf{$x < 2$} & 0,5đ \\
\hline
\textbf{Câu 16} & Giá trị biểu thức $N = \sqrt{(\sqrt{7}-3)^2} + \sqrt{7}$ & \textbf{$3$} & 0,5đ \\
\hline
\textbf{Câu 17} & Chiều cao tòa tháp ($15 \cdot \tan 55^\circ$, làm tròn 2 chữ số thập phân) & \textbf{$21{,}42\text{ m}$} & 0,5đ \\
\hline
\textbf{Câu 18} & Diện tích giấy làm quạt ($R = 20\text{ cm}, n = 120^\circ$) & \textbf{$418{,}7\text{ cm}^2$} & 0,5đ \\
\hline
\end{tabular}
\end{center}

\vspace{3mm}
\noindent
\textbf{\large PHẦN II. TỰ LUẬN (3,0 ĐIỂM)}

\vspace{2mm}
\textbf{Câu 19 (1,0 điểm).}
\begin{itemize}
  \item[a)] \textit{(0,5 điểm)} Giải hệ phương trình: $\begin{cases} 3x - y = 7 \\ 2x + 3y = 1 \end{cases}$\\
  Nhân phương trình thứ nhất với 3: $\begin{cases} 9x - 3y = 21 \\ 2x + 3y = 1 \end{cases}$ \hfill \textbf{(0,25đ)}\\
  Cộng hai phương trình: $11x = 22 \Leftrightarrow x = 2$. Thay vào phương trình đầu ta được: $y = 3(2) - 7 = -1$.\\
  Vậy hệ phương trình có nghiệm duy nhất: $(x; y) = (2; -1)$. \hfill \textbf{(0,25đ)}

  \item[b)] \textit{(0,5 điểm)} Rút gọn biểu thức $A$ với $x > 0, x \ne 9, x \ne 1$:\\
  Quy đồng mẫu thức trong ngoặc với mẫu chung là $(\sqrt{x}+3)(\sqrt{x}-3) = x - 9$:
  \[
  \dfrac{\sqrt{x}(\sqrt{x}-3) + 2\sqrt{x}(\sqrt{x}+3) - (3x+9)}{(\sqrt{x}+3)(\sqrt{x}-3)} = \dfrac{3\sqrt{x}-9}{(\sqrt{x}+3)(\sqrt{x}-3)} = \dfrac{3}{\sqrt{x}+3} \hfill \textbf{(0,25đ)}
  \]
  Chia cho phân thức thứ hai:
  \[
  A = \dfrac{3}{\sqrt{x}+3} : \dfrac{\sqrt{x}-1}{\sqrt{x}-3} = \dfrac{3(\sqrt{x}-3)}{(\sqrt{x}+3)(\sqrt{x}-1)} \hfill \textbf{(0,25đ)}
  \]
\end{itemize}

\vspace{1mm}
\textbf{Câu 20 (1,5 điểm).}
\begin{itemize}
  \item[a)] \textit{(0,75 điểm)} Vì $MA, MC$ là các tiếp tuyến của $(O)$ nên $MA \perp OA$ và $MC \perp OC \Rightarrow \widehat{MAO} = \widehat{MCO} = 90^\circ$. \hfill \textbf{(0,25đ)}\\
  Gọi $I$ là trung điểm của đoạn thẳng $MO$. Trong các tam giác vuông $MAO$ và $MCO$ có $AI, CI$ lần lượt là các đường trung tuyến ứng với cạnh huyền $MO \Rightarrow IA = IC = IM = IO = \dfrac{MO}{2}$. \hfill \textbf{(0,25đ)}\\
  Suy ra bốn điểm $A, M, C, O$ cùng thuộc đường tròn tâm $I$ đường kính $MO$. \hfill \textbf{(0,25đ)}

  \item[b)] \textit{(0,75 điểm)}
  \begin{itemize}
    \item $K \in (O)$ đường kính $AB \Rightarrow \widehat{AKB} = 90^\circ \Rightarrow AK \perp MB$. $\Delta MAB$ vuông tại $A$ có đường cao $AK \Rightarrow MA^2 = MK \cdot MB$. \hfill \textbf{(0,25đ)}
    \item Vì $CH \parallel MA$ (cùng vuông góc $AB$), trong $\Delta BAM$ có $NH \parallel MA \Rightarrow \dfrac{NH}{MA} = \dfrac{BH}{BA}$ (1). \hfill \textbf{(0,25đ)}
    \item Kéo dài $BC$ cắt tia $MA$ tại $E$. Ta có $OM \perp AC$ và $BC \perp AC \Rightarrow OM \parallel BE$. Trong $\Delta ABE$, $O$ là trung điểm $AB$ và $OM \parallel BE \Rightarrow M$ là trung điểm $AE$, tức $EA = 2MA$. Áp dụng định lí Thales: $\dfrac{CH}{EA} = \dfrac{BH}{BA}$ (2).\\
    Từ (1) và (2) suy ra $\dfrac{CH}{2MA} = \dfrac{NH}{MA} \Rightarrow CH = 2NH \Rightarrow N$ là trung điểm của $CH$. \hfill \textbf{(0,25đ)}
  \end{itemize}
\end{itemize}

\vspace{1mm}
\textbf{Câu 21 (0,5 điểm).}
\begin{itemize}
  \item[a)] \textit{(0,25 điểm)} Gọi cạnh đáy hình vuông là $x > 0$ (m) và chiều sâu là $h > 0$ (m).\\
  Thể tích $V = x^2 h = 32 \Rightarrow h = \dfrac{32}{x^2}$. Diện tích vật liệu xây dựng (bể không nắp):
  \[
  S = x^2 + 4xh = x^2 + \dfrac{128}{x} = x^2 + \dfrac{64}{x} + \dfrac{64}{x} \hfill \textbf{(0,10đ)}
  \]
  Áp dụng bất đẳng thức AM-GM cho 3 số dương: $S \ge 3\sqrt[3]{x^2 \cdot \dfrac{64}{x} \cdot \dfrac{64}{x}} = 3 \cdot 16 = 48\text{ m}^2$. \hfill \textbf{(0,075đ)}\\
  Dấu ``='' xảy ra khi $x^2 = \dfrac{64}{x} \Leftrightarrow x^3 = 64 \Leftrightarrow x = 4\text{ m} \Rightarrow h = 2\text{ m}$.\\
  \textbf{Kết luận:} Đáy hình vuông cạnh $4\text{ m}$, chiều sâu $2\text{ m}$, diện tích nhỏ nhất là $48\text{ m}^2$. \hfill \textbf{(0,075đ)}

  \item[b)] \textit{(0,25 điểm)} Với mọi số thực dương $t > 0$, áp dụng bất đẳng thức AM-GM:
  \[
  t^2 + \dfrac{2}{t} = t^2 + \dfrac{1}{t} + \dfrac{1}{t} \ge 3\sqrt[3]{t^2 \cdot \dfrac{1}{t} \cdot \dfrac{1}{t}} = 3 \hfill \textbf{(0,10đ)}
  \]
  Áp dụng lần lượt cho $x, y, z$ rồi cộng lại: $P = \sum \left(x^2 + \dfrac{2}{x}\right) \ge 3 + 3 + 3 = 9$. \hfill \textbf{(0,075đ)}\\
  Đẳng thức xảy ra khi $x = y = z = 1$. Khi đó $x + y + z = 3 \le 3$ (thỏa mãn giả thiết đề bài).\\
  \textbf{Kết luận:} Giá trị nhỏ nhất của $P$ là $9$, đạt được khi $x = y = z = 1$. \hfill \textbf{(0,075đ)}
\end{itemize}

\end{document}
